\documentclass[10pt,a4paper]{article}
\usepackage[a4paper,margin=1.8cm]{geometry}
\usepackage{multicol}
\usepackage{titlesec}
\usepackage{graphicx}
\usepackage{enumitem}
\usepackage{hyperref}
\usepackage{booktabs}
\usepackage{array}
\setlength{\columnsep}{1cm}
\titleformat{\section}{\bfseries\large}{\thesection.}{0.5em}{}
\titleformat{\subsection}{\bfseries}{\thesubsection}{0.5em}{}

\title{AI-driven Sustainable Meal Planning in Institutional Canteens}
\author{Samuele Mega}
\date{}

\begin{document}

\maketitle

\begin{multicols}{2}

\section{Project Idea}

This project introduces an AI-based system to intelligently plan weekly menus in institutional canteens (schools, hospitals, and companies). The primary goal is to balance sustainability, nutritional quality, and user satisfaction by:

\begin{itemize}[leftmargin=*]
  \item Gathering user preferences (both ranked recipes and portion sizes);
  \item Selecting recipes that maximize collective satisfaction;
  \item Choosing ingredients that optimize nutritional balance, environmental impact, and stock usage;
  \item Following dietary guidelines such as those from WHO and CREA.
\end{itemize}

Each recipe is defined as a flexible blueprint: categories (e.g. \emph{cereals, legumes, dairy}) with relative proportions, rather than fixed ingredients.

The system is highly modular. Preferences, recipes, warehouse data, and nutritional / eco properties are stored separately, allowing reuse and adaptation to different institutions or dietary rules. This decoupling supports extensibility: for example, it would be simple to plug in additional modules for special diets (gluten-free, vegan, halal, etc.) or policy constraints (budget, calories).

\begin{center}
\begin{tabular}{@{} l l l @{} }
\toprule
\multicolumn{3}{c}{\textbf{Recipe blueprint example}} \\
\toprule
\multicolumn{3}{c}{\textbf{Ingredients}} \\
\midrule
\textbf{Category} & \textbf{Alternatives} & \textbf{Qty (\%)} \\
\midrule
Cereals & Rice, Barley, Spelt & 50 \\
Legumes & Chickpeas, Lentils & 30 \\
Dairy & Ricotta, Parmigiano & 15 \\
Oil & Olive Oil, Avocado Oil & 5 \\
\midrule
\multicolumn{3}{c}{\textbf{Macronutrient Distribution Targets}} \\
\midrule
\textbf{Nutrient} & \textbf{Target (\%)} & \\
Carbohydrates & 60 & \\
Proteins & 20 & \\
Fats & 20 & \\
\midrule
\multicolumn{3}{c}{\textbf{Allowed Portion Sizes}} \\
\midrule
\textbf{Size} & \textbf{Weight (g)} & \\
Small & 300 & \\
Medium & 400 & \\
Large & 500 & \\
\bottomrule
\end{tabular}
\end{center}

These structured constraints are directly encoded in the chromosome and evaluation logic of the genetic algorithm, which guide the optimization process to produce feasible and nutritionally sound meal configurations.

\section{Social Impact}

In Italy, institutional food services waste approximately 30\% of prepared food annually, contributing to more than \euro14 billion in national food waste. This project addresses inefficiencies in demand forecasting, menu design, and inventory management by:

\begin{itemize}[leftmargin=*]
  \item Aligning menus with what people actually want to eat, reducing uneaten portions;
  \item Making use of existing stock to minimize unnecessary procurement and reduce spoilage;
  \item Educating users through diverse meals that reflect seasonal, sustainable choices;
  \item Providing an open-source, reproducible framework that can be customized and scaled.
\end{itemize}

This approach promotes responsible consumption and healthy diets, especially for children and vulnerable populations in public facilities. In addition, it could influence public procurement strategies and sustainability certifications.

\section{AI Techniques}

\subsection{Recipe Selection via Ranked Pairs}

Each user submits a ranking on the available recipes. To fairly aggregate preferences, we apply the \textbf{ ranked pairs} method from the voting theory. It is a \textit{Condorcet-consistent} algorithm that selects the widely preferred options without falling into majority cycles. Concretely, the algorithm:

\begin{itemize}
  \item Computes pairwise victories and their margins from user ballots;
  \item Sorts the pairs by strength of victory;
  \item Iteratively "locks in" the strongest pairwise wins, as long as they don't introduce cycles in the preference graph.
\end{itemize}

We use \texttt{pyrankvote} Python library for this step. The result is a sorted list of recipes that best reflects collective preference, satisfying more users, and reducing the rejection of meals.

\subsection{Ingredient Optimization via Genetic Algorithm}

Once a recipe blueprint is selected (e.g., \emph{cereals + legumes + dairy + oil}), the system fills each ingredient slot using a \textbf{Genetic Algorithm (GA)}. This is a population-based metaheuristic for combinatorial optimization. The GA components are:
\begin{itemize}[leftmargin=*]
  \item \textbf{Chromosome}: a full assignment of foods to ingredient categories;
  \item \textbf{Gene}: a specific food (e.g., \emph{lentils}) and a quantity scaling factor (e.g., \emph{1.05});
  \item \textbf{Population}: initialized with random valid food combinations from the database;
  \item \textbf{Fitness function}: evaluates candidates by computing a weighted score from:
    \begin{itemize}[leftmargin=*]
      \item Deviation from macronutrient targets (carbs, proteins, fats);
      \item Eco-score and seasonality of ingredients;
      \item Availability in the warehouse (including penalties for missing stock);
      \item Price on the market.
    \end{itemize}
  \item \textbf{Operators}:
    \begin{itemize}[leftmargin=*]
      \item Crossover: mix ingredients from two parents;
      \item Mutation: randomly replace or rescale an ingredient.
    \end{itemize}
\end{itemize}

The GA runs for a fixed number of generations or until convergence, returning one or more optimized menus. This technique is highly flexible and allows real-world constraints to be seamlessly incorporated.

\subsection{Example Output}

Below is a sample output of the system for one selected recipe (e.g., \textit{Cereals and Legumes}):

\begin{center}
\begin{tabular}{@{} l l l l l l @{}} 
\toprule
\textbf{Food} & \textbf{W (kg)} & \textbf{C} & \textbf{P} & \textbf{F} & \textbf{Cost (€)} \\
\midrule
Spelt & 40.0 & 58.0 & 12.0 & 2.0 & 45.0 \\
Lentils & 24.0 & 20.0 & 9.0 & 1.0 & 30.0 \\
Ricotta & 12.0 & 3.0 & 6.0 & 9.0 & 22.0 \\
Olive oil & 4.0 & 0.0 & 0.0 & 100.0 & 10.0 \\
\midrule
\textbf{Total} & \textbf{80.0} & \textbf{39.0} & \textbf{13.0} & \textbf{6.75} & \textbf{107.0} \\
\bottomrule
\end{tabular}
\end{center}

The system generates multiple such combinations. Each is scored by a fitness function that balances macronutrient compliance, ecological impact, seasonality, cost, and stock availability.

\section{Links to AI and Other Courses}

This project integrates multiple concepts from AI and Computer Engineering curricula:
\begin{itemize}[leftmargin=*]
  \item \textbf{Artificial Intelligence}: metaheuristics (GA), voting systems, preference aggregation and selection with soft constraints;
  \item \textbf{Graph Algorithms}: cycle detection in Ranked Pairs (via DFS);
  \item \textbf{Multi-objective Optimization}: trade-offs between nutrition, ecology, and cost;
  \item \textbf{Data Engineering}: preprocessing and managing structured food and recipe datasets;
  \item \textbf{Software Design}: modular architecture, extensibility, reproducibility.
\end{itemize}

\section{Further Implementations}

This system could be extended in several meaningful directions:

\begin{itemize}[leftmargin=*]
  \item \textbf{User Feedback Loop}: Collect feedback on meals actually consumed and integrate it to improve recipe selection over time.
  \item \textbf{Dietary Constraints Integration}: Incorporate user-specific constraints such as allergies, intolerances, religious diets, or chronic conditions.
  \item \textbf{Reinforcement Learning}: Use RL agents to adapt the menu policy based on long-term performance in real deployments.
  \item \textbf{Web Interface \& API}: Provide a full-stack application where canteen managers can configure parameters and run AI through a browser.
\end{itemize}

\section*{GitHub Repository}
Code and documentation: \\ \url{https://github.com/samuelemega/unipd-ai-canteens}

\section*{Usage Instructions}
To run the project locally, use the following commands:

\begin{itemize}[leftmargin=*]
  \item \texttt{python3 scripts/random\_preferences.py}
  \item \texttt{python3 scripts/choose\_recipes.py}
  \item \texttt{python3 scripts/choose\_ingredients.py}
\end{itemize}

\end{multicols}
\end{document}
